\section{Introduction}
\label{sec:intro}

Representation models have received considerable attention in computer vision~\cite{beit,mae,ibot,moco,mocov2,mocov3,simclr,dino,dinov2}, speech processing~\cite{wav2vec,wav2vec2,wavlm,hubert}, natural language processing~\cite{elmo,bert,roberta,electra,deberta}, etc. 
Learning from large amounts of data, representation models demonstrate strong generalization ability in a wide range of downstream tasks.
Furthermore, the explosive growth of large-scale language models (LLMs) has sparked an escalating appetite for representation models.
% Representation models that learn from large amounts of data demonstrate strong generalization ability in downstream tasks, while the rapid development of large-scale language models (LLMs) has also created an increasing demand for representation models~\cite{gpt4,kosmos,blip2,minigpt4,llava,mplug-owl}.
% Recent research shows that integrating LLMs with representation models enables LLMs to understand other modalities (e.g., image), and strong representation models can lead to better performance~\cite{blip2}.
Until recently, representation models have shown their bedrock role to unleash LLMs to understand, perceive, and interact with other modalities (e.g., vision)~\cite{gpt4,kosmos,blip2,minigpt4,llava,mplug-owl,palme}.

Due to the distinct characteristics of different modalities, previous research mainly focuses on building uni-modal representation models with individual architectures and pretraining tasks.
Despite achieving excellent results, uni-modal representation models face difficulties in effectively utilizing multi-modal data such as image-text pairs and audio-text pairs, which makes them challenging to extend to multi-modal tasks.
% In addition, humans generally acquire knowledge about the world by merging information from multiple modalities, which is more effective than learning from single modalities.
% Some studies have also shown that introducing multi-modal information can improve model performance in uni-modal tasks~\cite{ofa,icode,fd-clip,eva}.
With the development of unified architectures~\cite{transformer,vit,perceiver,perceiverio} and efficient pretraining tasks~\cite{bert,beit,mae,clip}, recent works have achieved promising results in vision-language learning~\cite{ofa,flamingo,coca,beit3,blip,pali} and audio-language learning~\cite{audioclip,wav2clip,clap,speecht5}. Nevertheless, there is still rare research on developing general models that can be applied to vision, audio, and language modalities.
\cite{beit3} utilize the Multiway Transformer to process both image and text modalities with a unified masked prediction task for pretraining. The masked prediction task requires a pretrained CLIP~\cite{clip} model to discretize image data, which limits the scalability to other modalities such as audio.
\cite{data2vec} proposes a general pretraining method that can be applied to vision, audio, and language modalities without the need for third-party models (e.g., CLIP), but it doesn't extend the method to multi-modal data.
% , therefore cannot handle multi-modal tasks.

In this paper, we explore a scalable way to build a general representation model toward unlimited modalities.
We advocate that a general representation model should meet the following conditions:
1. The model architecture must be flexible enough to accommodate various modalities and support multi-modal interaction. 
2. Pretraining tasks should not only extract information within each modality but also ensure alignment across modalities. 
3. Pretraining tasks should be general and straightforward, allowing them to be applied to different modalities.

% In this paper, we explore a scalable way to build a general representation model toward unlimited modalities.
% We advocate that a general representation model should meet the following conditions: 
% 1). The model architecture needs to be sufficiently flexible. Ideally, the model should incorporate modality-specific components to capture modality-specific information, as well as modality-interaction components to enable cross-modal interaction. 
% 2). The pretraining tasks should not only focus on extracting information within each modality but also ensure alignment across modalities, so that the model can achieve superior performance in both uni-modal and multi-modal tasks.
% 3). The pretraining tasks should be general and simple enough to be applied to different modalities.
% There is no need to discretize data using external models, nor is it necessary to make specific designs for each modality.

Driven by these motivations, we propose \onepeace, a model with 4B parameters that can seamlessly align and integrate representations across vision, audio, and language modalities.
The architecture of \onepeace consists of multiple modality adapters and a modality fusion encoder.
% The modality fusion encoder is based on the Transformer architecture. Each Transformer block contains a shared self-attention layer and multiple modality-specific Feed Forward Networks (FFN).
% The modality-specific adapters are responsible for converting raw inputs of images, audio, and text to features that can be processed by the modality fusion encoder. 
Each modality is equipped with an adapter for converting the raw inputs into feature sequences.
The modality fusion encoder operates on feature sequences with Transformer architecture. Each Transformer block contains a shared self-attention layer and multiple modality Feed Forward Networks (FFNs).
The self-attention layer enables interaction between the multi-modal features through the attention mechanism, while the modality FFNs facilitate information extraction within modalities.
With the clear division of labor in this architecture, extending new modalities only requires the injection of adapters and FFNs.

To pretrain \onepeace, we design two modality-agnostic pretraining tasks.
The first one is cross-modal contrastive learning, it contains both vision-language contrastive learning and audio-language contrastive learning, which effectively align the semantic spaces of vision, audio, and language modalities.
The second one is intra-modal denoising contrastive learning, it can be viewed as a combination of masked prediction~\cite{bert,beit,mae,wav2vec} and contrastive learning~\cite{clip,simclr,moco}, where we perform contrastive loss between the fine-grained masked features and visible features, such as image patches, language tokens, or audio waveform features.
These tasks collaborate to enhance the model's fine-tuning performance while also maintaining cross-modal retrieval capability.
Furthermore, they are universal for all modalities, which obviates the need for modality-specific designs.
With the scaling-friendly model architecture and pretraining tasks, \onepeace has the potential to expand to unlimited modalities.

We conduct comprehensive experiments on different tasks across various modalities, including vision, audio, vision-language, and audio-language tasks.
Without using any vision or language pretrained model for initialization, \onepeace achieves leading results in both uni-modal and multi-modal tasks, including image classification ($89.8\%$ accuracy on ImageNet w/o privately labeled data), semantic segmentation ($63.0\%$ mIoU on ADE20K), audio-text retrieval (outperforming previous SOTAs on AudioCaps and Clotho by a large margin), audio classification ($91.8\%$ zero-shot accuracy on ESC-50, $69.7\%$ accuracy on FSD50K, $59.6\%$ accuracy on VGGSound w/o visual information), audio question answering ($86.2\%$ accuracy on AVQA w/o visual information), image-text retrieval ($84.1\%$ I2T R@1 on MSCOCO and $97.6\%$ I2T R@1 on Flickr30K w/o intermediate finetuning and ranking), and visual grounding ($89.26\%$/$83.23\%$/$89.27\%$ scores on RefCOCO/+/g test sets).
% Also, after fine-tuning on the visual grounding dataset, \onepeace is capable of transferring its visual grounding ability to out-of-domain pictures.

% To summarize:

% \begin{itemize}
%     \item We propose \onepeace, a general representation model with 4B parameters that can seamlessly align and integrate representations across vision, audio, and language modalities
%     With the scaling-friendly model architecture and pretraining tasks, \onepeace has the potential to expand to infinite modalities.
%     \item \onepeace, as a general representation model, achieves state-of-the-art results in uni-modal tasks, including image classification (ImageNet), semantic segmentation (ADE20K) and audio classification (ESC-50, FSD50K, VGGSound).
%     \item Without using vision or language pretrained models for initialization, \onepeace achieves state-of-the-art results in multi-modal tasks, including audio-text retrieval (AudioCaps, Clotho), audio question answering (AVQA), image-text retrieval (MSCOCO, Flickr30K), and visual grounding (RefCOCO/+/g).
% \end{itemize}


